\documentclass[11pt]{article}
\usepackage[T1]{fontenc}
\usepackage{lmodern,amsmath,amssymb,amsthm,booktabs,tabularx}
\usepackage[margin=1in]{geometry}
\usepackage{xurl}
\usepackage[colorlinks=true,linkcolor=blue,urlcolor=blue,citecolor=blue]{hyperref}
\hypersetup{pdftitle={Factorial bounds and sharp onsets for A281428: extending Su (2026) on Bala's COMMENTS, items 1–2},
pdfauthor={Haobo Ma; Wenlin Zhang},pdfsubject={Bala's COMMENTS items 1–2; CC BY 4.0}}
\newtheorem{theorem}{Theorem}
\newtheorem{lemma}[theorem]{Lemma}
\newtheorem{corollary}[theorem]{Corollary}
\newcommand{\Q}{\mathbb{Q}}
\newcommand{\Z}{\mathbb{Z}}
\newcommand{\coeff}[2]{[#1]#2}
\title{Factorial bounds and sharp onsets for A281428:\\
extending Su (2026) on Bala's \textsc{Comments}, items 1--2}
\author{Haobo Ma\\\small ChronoAI Pte Ltd, Singapore
\and Wenlin Zhang\\\small National University of Singapore\\\small ChronoAI Pte Ltd, Singapore}
\date{November 11, 2026}
\begin{document}
\maketitle
\begin{abstract}
We extend Su's (2026) proof of Peter Bala's conjectures in OEIS A281428,
\textsc{Comments}, items 1--2 (displayed March 3, 2026), with factorial bounds and
sharp onsets for the two accompanying examples.
For its coefficients $a(n)=(2n)!\coeff{x^{2n}}{C(x)}$, we prove
$v_p(a(n))\geq v_p((2n)!)$ for every prime $p\geq5$.
The proof combines integral exponential coefficients with ordinary coefficients
in $\Z[1/6]$, using the inverse cubic $u-u^3/3=x$.
The exact eventual-zero onset modulo $385$ is $6$; modulo $27$, the least
eventual period is $18$ and its earliest onset is $2$.
The same argument gives factorial bounds for every polynomial in the two
differential-system solutions, including the companion A281427 and A052502.
\end{abstract}

\section{The posed conjecture and a valuation bound}
The entry A281428, introduced by Paul D. Hanna, defines $C\in\Q[[x]]$ by
\begin{equation}\label{eq:source}
C(x)=\cosh\left(\int_0^x C(t)^4\,dt\right),\qquad
C(0)=1,\qquad a(n)=(2n)!\coeff{x^{2n}}{C(x)}.
\end{equation}
All series and integrals are formal. In \textsc{Comments}, item 1,
Peter Bala conjectured eventual divisibility of $a(n)$ by every positive $N$
coprime to $6$ \cite{oeis}. Item 2 asserts that when $2\mid N$ or $3\mid N$,
the sequence modulo $N$ has an eventual period dividing $\varphi(N)$,
where $\varphi$ is Euler's totient function.
Both items carry Bala's signature of March 3, 2026.
The edit history records their entry on February 16 (revision 32),
the clarification of item 2 on February 17 (revision 33), and the addition
of the examples and displayed March date in revision 37.
Su \cite[Theorem 4.1, pp.~7--8]{su} proved
this assertion and the neighboring eventual-period conjecture.
Here the inverse cubic already recorded in the entry yields a direct
denominator argument and the following factorial bound.
For a nonzero rational number $q$, $v_p(q)$ is its $p$-adic valuation;
we put $v_p(0)=+\infty$.

\begin{theorem}\label{thm:main}
Equation \eqref{eq:source} has a unique solution, which is even.
Its coefficients satisfy $a(n)\in\Z$ and
\begin{equation}\label{eq:valuation}
v_p(a(n))\ \geq\ v_p((2n)!)\qquad(n\geq0,\ p\geq5\text{ prime}).
\end{equation}
Consequently, if $(N,6)=1$ and $N\mid(2n)!$, then $N\mid a(n)$.
In particular $n\geq\lceil N/2\rceil$ is sufficient.
\end{theorem}

\begin{proof}
Introduce $S=\sinh(\int_0^x C^4)$.
The defining equation gives the differential system
\begin{equation}\label{eq:ode}
S'=C^5,\qquad C'=SC^4,\qquad S(0)=0,\quad C(0)=1.
\end{equation}
Conversely, the coefficients of this system are determined successively
over $\Q$: at degree $j+1$ the right sides use coefficients through degree $j$.
Thus it has a unique solution. Its invariant is $C^2-S^2=1$, since
the derivative of the left side is zero and its constant coefficient is one.
Writing $T=\int_0^x C^4$, the equations
$(C+S)'=T'(C+S)$ and $(C-S)'=-T'(C-S)$ with constant coefficient one give
$C+S=e^T$ and $C-S=e^{-T}$. This recovers \eqref{eq:source}.
Uniqueness applied to $(-S(-x),C(-x))$ proves the asserted parity.

To prove integrality with factorial normalization, work in
\[
R=\Z[s,c]/(c^2-s^2-1),\qquad
D(s)=c^5,\quad D(c)=sc^4,\quad E(f)=f(0,1).
\]
The derivation is well defined because
$D(c^2-s^2-1)=2sc^5-2sc^5=0$.
The map $f\mapsto\sum_{j\geq0}E(D^j f)x^j/j!$ is multiplicative:
the coefficient identity is the iterated Leibniz rule
$D^j(fg)=\sum_{i=0}^j\binom jiD^i(f)D^{j-i}(g)$.
Applied to $s,c$ it therefore solves \eqref{eq:ode}.
In particular
\begin{equation}\label{eq:integral}
j!\coeff{x^j}{C}=E(D^jc)\in\Z.
\end{equation}

Set $u=S/C$, which is defined since $C(0)=1$. The invariant and
\eqref{eq:ode} imply
\begin{equation}\label{eq:cubic}
u'=C^2,\qquad (1-u^2)C^2=1,\qquad
u-\frac{u^3}{3}=x.
\end{equation}
For the last equality, differentiate: $(1-u^2)u'=1$; both sides have
constant coefficient zero before differentiation.
The recurrence $u=x+u^3/3$, with $u(0)=0$, shows by degree induction that
$u\in\Z[1/3][[x]]$. Indeed the coefficient of degree $j$ in $u^3$
uses coefficients of $u$ of degrees strictly less than $j$.
It follows that $C^2=(1-u^2)^{-1}\in\Z[1/3][[x]]$.
If $C=1+\sum_{j\geq1}c_jx^j$ and $C^2=1+\sum_{j\geq1}b_jx^j$, then
\[
2c_j=b_j-\sum_{i=1}^{j-1}c_i c_{j-i}.
\]
Another degree induction gives $C\in\Z[1/6][[x]]$.
Hence $v_p(\coeff{x^{2n}}{C})\geq0$ for $p\geq5$.
Combining this with \eqref{eq:integral} proves \eqref{eq:valuation}.
For $(N,6)=1$, apply the valuation inequality at every prime dividing $N$.
Finally $2n\geq N$ ensures $N\mid(2n)!$.
\end{proof}

\begin{corollary}[Polynomial observables]\label{cor:polynomial}
For every $F\in\Z[s,c]$ and $j\geq0$, the number
$b_F(j)=j!\coeff{x^j}{F(S(x),C(x))}$ is an integer, and
$v_p(b_F(j))\geq v_p(j!)$ for every prime $p\geq5$.
In particular this holds for all columns
$A(n,k)=(2n)!\coeff{x^{2n}}{C^k}$ with $k\geq0$.
\end{corollary}
\begin{proof}
The multiplicative differential-series map in the proof gives
$b_F(j)=E(D^jF)\in\Z$. Since $S=uC$, both $S$ and $C$ lie in
$\Z[1/6][[x]]$, as does $F(S,C)$. Multiplying its $p$-integral
coefficient by $j!$ proves the bound.
\end{proof}

\begin{corollary}\label{cor:threshold}
For $N\geq1$ coprime to $6$, let
\[
\mu(N)=\min\{m\geq0:N\mid m!\}.
\]
Then $N\mid a(n)$ for all $n\geq\lceil\mu(N)/2\rceil$.
For $N=385$ the least eventual-zero onset is $6$.
\end{corollary}
\begin{proof}
The first assertion follows from Theorem~\ref{thm:main}.
Since $385=5\cdot7\cdot11$ divides $11!$ but not $10!$, $\mu(385)=11$.
The initial values give $a(5)=15335425\equiv105\pmod{385}$,
so an onset at most $5$ is impossible. This proves Bala's
\textsc{Example} about divisibility modulo $385$ for every $n\geq6$.
The same onset is already a consequence of Su's Theorem 4.1(1), whose
$r=1$ bounds for $p=5,7,11$ are $3,4,6$, together with the displayed $a(5)$.
\end{proof}
For a prime power $p^r$ with $p\geq5$, a convenient explicit sufficient onset is
$\lceil rp/2\rceil$, because $v_p((2n)!)\geq\lfloor2n/p\rfloor$.
The sharper factorial criterion uses
$v_p(m!)=\sum_{j\geq1}\lfloor m/p^j\rfloor$.

\section{The excluded primes and source coverage}
The same differential system gives a short converse. Following the
parameterized recurrence credited to Seiichi Manyama in A281428
(April 14, 2026), write
$A(n,k)=(2n)!\coeff{x^{2n}}{C^k}$ for integers $k\geq0$.
\begin{lemma}\label{lem:columns}
For $k\geq0$, $A(0,k)=1$, and for $n\geq1$,
\begin{equation}\label{eq:manyama}
A(n,k)=k(k+4)A(n-1,k+8)-k(k+3)A(n-1,k+6).
\end{equation}
Every $a(n)$ is odd, and $a(n)\equiv(-1)^{n-1}\pmod3$ for $n\geq1$.
\end{lemma}
\begin{proof}
Differentiate $C^k$ twice and use $S^2=C^2-1$:
\[
(C^k)''=kC^{k+8}+k(k+3)S^2C^{k+6}
=k(k+4)C^{k+8}-k(k+3)C^{k+6}.
\]
Taking the $(2n-2)$nd derivative at zero proves \eqref{eq:manyama}.
For odd $k$, its first multiplier is odd and its second is even;
induction proves $A(n,k)$ odd.
Modulo $3$ we prove $A(n,k)\equiv(-1)^{n-1}k$ for every $k$ and $n\geq1$.
The case $n=1$ is $A(1,k)=k$. For $n\geq2$, substitute the induction
hypothesis in \eqref{eq:manyama}; apart from the factor $(-1)^{n-2}$,
the resulting expression is
\[
k(k+1)(k+2)-k^3\equiv-k\pmod3.
\]
Here three consecutive integers have product divisible by $3$, and
$k^3\equiv k\pmod3$. Set $k=1$.
\end{proof}
\begin{corollary}
For $N\geq1$, $N$ divides all sufficiently late $a(n)$ if and only if
$(N,6)=1$. In fact no $a(n)$ is divisible by $2$ or $3$.
\end{corollary}
\begin{proof}
Use Theorem~\ref{thm:main} and Lemma~\ref{lem:columns}; $a(0)=1$.
\end{proof}

\begin{corollary}[Two cross-referenced companions]\label{cor:companions}
With the original offsets, put
$h(n)=(2n-1)!\coeff{x^{2n-1}}S$ for $n\geq1$ (A281427), and
$r(n)=(2n)!\coeff{x^{2n}}{C^2}$ for $n\geq0$ (A052502).
For $p\geq5$, $v_p(h(n))\geq v_p((2n-1)!)$.
A positive $N$ eventually divides $h(n)$ exactly when $(N,6)=1$;
a sufficient onset is $\lceil(\mu(N)+1)/2\rceil$.
Moreover
\[
r(n)=\frac{(3n)!}{3^n n!},\qquad
v_p(r(n))\geq v_p((2n)!)\quad(p\ne3).
\]
A positive $N$ eventually divides $r(n)$ exactly when $3\nmid N$,
with sufficient onset $\lceil\mu(N)/2\rceil$.
\end{corollary}
\begin{proof}
The first bound follows from Corollary~\ref{cor:polynomial}.
Since $S'=C^5$, $h(n)=A(n-1,5)$; Lemma~\ref{lem:columns} shows these
numbers are odd and $h(n)\equiv(-1)^{n-1}\pmod3$ (including $n=1$).
The stated factorial threshold then proves the characterization.
Lagrange inversion, as displayed in Section~\ref{sec:coverage}, gives
$\coeff{x^{2n}}{C^2}=3^{-n}\binom{3n}{n}$ and the formula for $r(n)$.
This coefficient is $p$-integral for every $p\ne3$.
Also $r(n)=\prod_{j=1}^n(3j-1)(3j-2)$ is an integer prime to $3$.
These facts prove the second characterization and threshold.
\end{proof}
These are consequences for the named companions, not conjectures posed
in their entries. Their source identities and offsets are given in
\cite{companions}; no assertion of priority is made for these consequences.

\section{The modulo-27 example: exact period and onset}\label{sec:27}
\begin{theorem}\label{thm:27}
Modulo $27$, $a(n)$ has least eventual period $18$, and the earliest
onset for this period is $n=2$.
\end{theorem}
\begin{proof}
We use the recurrence credited to Kotesovec in A281428 (February 24, 2026):
\begin{equation}\label{eq:kotesovec}
\begin{split}
16a(n)={}&(288n^2-576n+317)a(n-1)\\
&-9(2n-3)^2(6n-11)(6n-7)a(n-2),\quad n\geq2.
\end{split}
\end{equation}
For completeness, it follows directly from the inverse cubic.
Writing $u=S/C$, differentiation gives
\begin{align*}
C'&=u(1-u^2)^{-5/2},&
C''&=(1+4u^2)(1-u^2)^{-9/2},\\
C'''&=u(17+28u^2)(1-u^2)^{-13/2}.
\end{align*}
Substitution of these expressions and $x=u-u^3/3$ yields
\begin{equation}\label{eq:linear}
(9x^2-4)^2C'''+54x(9x^2-4)C''+(531x^2-101)C'+45xC=0.
\end{equation}
Indeed, after factoring out $u(1-u^2)^{-5/2}$, the remaining
polynomial, with $t=u^2$, is
\begin{align*}
&(17+28t)(4-t)^2-18(3-t)(4-t)(1+4t)\\
&\hspace{2em}+59t(3-t)^2-101+15(3-t)(1-t)^2=0.
\end{align*}
The coefficient of $x^{2n-3}/(2n-3)!$ in \eqref{eq:linear} is
\eqref{eq:kotesovec} for every $n\geq2$.

Reducing that recurrence modulo $27$ gives
\begin{equation}\label{eq:mod27}
16a(n)\equiv(18n^2+18n+20)a(n-1)+9n^2a(n-2)\pmod{27}.
\end{equation}
Its coefficients have period $3$, and $16$ is invertible modulo $27$.
Starting from $a(0)=a(1)=1$, the exact residue calculation gives
\[
\begin{array}{c|rrrrrrrrrr}
n&2&3&4&5&6&7&8&9&10&11\\ \hline
a(n)\bmod27&17&1&8&19&17&1&26&19&17&10
\end{array}
\]
\[
\begin{array}{c|rrrrrrrrrr}
n&12&13&14&15&16&17&18&19&20&21\\ \hline
a(n)\bmod27&26&19&8&10&26&1&8&10&17&1.
\end{array}
\]
Thus the pairs at $2,3$ and $20,21$ agree. Induction in
\eqref{eq:mod27} proves $a(n+18)\equiv a(n)$ for every $n\geq2$.
An earlier onset is impossible because $a(1)=1\not\equiv10=a(19)$.
On a sequence with period $18$, its least eventual period divides $18$:
any tail period also acts on the same $18$ residue classes.
Every proper divisor of $18$ divides $6$ or $9$.
But $a(2)=17\not\equiv26=a(8)$ and
$a(2)=17\not\equiv10=a(11)\pmod{27}$, excluding both periods.
\end{proof}
This proves the source's numerical observation beginning at $7$, and
sharpens its onset to $2$. Su's Theorem 4.1(3) already gives period $18$
from $3$; the present recurrence argument identifies the optimal onset
and the least period for this particular modulus.
The optimality conclusions also follow by combining Su's theorem with
Bala's displayed residues; the proof above does not require the general
period theorem.

\begin{corollary}[A cross-referenced operator identity]\label{cor:operator}
In A281183, \textsc{Formula}, December 22, 2025, define polynomials
$g_0(z)=1$ and $g_{j+1}(z)=(1+z^2)\frac{d}{dz}((1+z^2)g_j(z))$.
Then $g_{2n}(0)$ is the $n$th term of A281183, as suggested there.
\end{corollary}
\begin{proof}
Write $M(\theta)=\sec^2\theta$. Substitution $z=\tan\theta$ gives
$g_{j+1}(\tan\theta)=\frac{d}{d\theta}(M(\theta)g_j(\tan\theta))$.
Let $V(0)=0$, $V'=M(V)$; thus $V$ is the inverse of
$\int_0^\theta\cos^2 t\,dt$, and $V'$ is the source e.g.f. of A281183.
The operator identity
$(\partial_\theta M)^j=M^{-1}(M\partial_\theta)^jM$ yields
$g_j(0)=\left.\frac{d^j}{dx^j}M(V(x))\right|_{x=0}$, since $M(0)=1$.
For $j=2n$ this is exactly the source factorial normalization.
\end{proof}

\section{Source coverage and the Lean declaration}\label{sec:coverage}
The two conjectures and two numerical observations in Bala's
March 3, 2026 contribution have the following coverage.
\begin{center}\small
\begin{tabularx}{\linewidth}{@{}p{.28\linewidth}X@{}}
\toprule
Source item in A281428 & Mathematical coverage\\
\midrule
\textsc{Comments}, 1 & Reproved here with a factorial valuation bound;
prior proof: Su (2026), Theorem 4.1(1).\\
\textsc{Comments}, 2 & Prior proof cited: Su (2026), Theorem 4.1(2)--(3)
and its Chinese remainder argument. The eventual period divides $\varphi(N)$.\\
\textsc{Example}, modulus $385$ & Proved for all $n\geq6$ in
Corollary~\ref{cor:threshold}, including minimality of the onset.\\
\textsc{Example}, modulus $27$ & Reproved and sharpened in
Theorem~\ref{thm:27}: least period $18$, earliest onset $2$;
prior proof: Su (2026), Theorem 4.1(3), onset $3$.\\
Manyama's columns, \textsc{Formula}, April 14, 2026 & Proved recurrence
in Lemma~\ref{lem:columns}; all-column valuation bound in
Corollary~\ref{cor:polynomial} (not posed as a conjecture).\\
A281427 and A052502, cross-references & Not posed; proved consequences
in Corollary~\ref{cor:companions}.\\
\textsc{Formula}, January 22, 2026 & Refuted-and-corrected ordinary
coefficient normalization, following Su (2026), Remark 4.2; proof below.\\
A281180, A281181, A281183, cross-references & Prior proofs cited:
Su (2026), Corollary 3.3 and Propositions 3.4, 3.6; these use exponent $4$,
not the exponent-$5$ flow here.\\
A281183, \textsc{Formula}, December 22, 2025, $g$ identity & Proved in
Corollary~\ref{cor:operator}; a distinct operator identity, not a new
periodicity claim.\\
\bottomrule
\end{tabularx}
\end{center}
The inverse-cubic and differential identities used above are source formulas,
not additional conjectures. The other inverse-function representations and
the stated asymptotic are not covered here. The source's algebraic formula
$(4-9x^2)C^6-3C^2-1=0$ follows by putting $z=C^2$ in
$9x^2z^3=(z-1)(2z+1)^2$; the latter is the square of
$x=u-u^3/3$ with $u^2=(z-1)/z$.
For the source's operator $f_k(n,0)$ in the January 22, 2026
\textsc{Formula} block, Su's Remark 4.2 gives the corrected normalization
$f_k(n,0)=n!\coeff{x^n}{C(x)^k}$. Already $k=1,n=2$ distinguishes
$f_1(2,0)=1$ from $\coeff{x^2}{C}=1/2$.
To see the full identity, put $u=\sin\theta$, so that
$x=\sin\theta-\sin^3\theta/3$, $C=\sec\theta$, and
$d/dx=L=\sec^3\theta\,d/d\theta$.
The source recurrence for $f_k$ gives
$\sec^3\theta\,f_k(n,\theta)=L^n\sec^k\theta$ by induction from
$f_k(0,\theta)=\sec^{k-3}\theta$.
Evaluation at $\theta=0$ proves the factorial correction.

For a direct rational-coefficient check, Lagrange inversion
\cite[Theorem 2.1.1, equation (2.1.1)]{gessel} applied to
$u=x/(1-u^2/3)$ gives
\[
\coeff{x^{2n+1}}{u}
=\frac1{2n+1}\coeff{t^{2n}}{(1-t^2/3)^{-(2n+1)}}
=\frac1{(2n+1)3^n}\binom{3n}{n}.
\]
Since $C^2=u'$, this recovers the source identity
$\coeff{x^{2n}}{C^2}=3^{-n}\binom{3n}{n}$.
The accompanying \texttt{verify.py} uses this identity and square-root
coefficient recursion through $n=128$. It compares two integer recurrences
through $n=256$, checks the displayed initial values, and checks
the valuation inequalities for every prime $5\leq p\leq101$ over that range.
These are finite checks; the all-index valuation bound is proved above.

The pinned Lean file
\href{https://github.com/the-omega-institute/trureturing/blob/b0f7c01b01511ea00428c1b6edad5dd2b372eaf8/D5/S1/Recurrence/Residue/BalaQuarticEGFEventualDivisibility.lean}{\texttt{BalaQuarticEGFEventualDivisibility.lean}}
contains the declaration \texttt{result} in namespace
\texttt{D5.S1.Recurrence.Residue.}\allowbreak
\texttt{BalaQuarticEGFEventualDivisibility}.
Its statement is existence of a solution to \eqref{eq:source} and eventual
integer divisibility for every positive $N$ coprime to $6$;
the proof uses the sufficient onset $\lfloor(N+1)/2\rfloor$.
The factorial valuation bound, the converse at $2$ and $3$, and the sharp
onset modulo $385$, polynomial-observable and companion corollaries,
the cross-referenced operator identity, and exact modulo-$27$ period and
onset have written proofs here.
These additions are not statements of that declaration.
Su's cited theorem supplies
the general eventual-period statement.

\begin{thebibliography}{9}
\bibitem{oeis}
The OEIS Foundation Inc., \emph{A281428}.
Entry by Paul D. Hanna (January 21, 2017);
conjectures by Peter Bala, \textsc{Comments}, items 1--2
(displayed March 3, 2026; first entered February 16, 2026).
\url{https://oeis.org/A281428}. Accessed November 11, 2026.
\bibitem{companions}
The OEIS Foundation Inc., \emph{A281427} and \emph{A052502}.
\url{https://oeis.org/A281427}; \url{https://oeis.org/A052502}.
Companion identities credited to Peter Bala (December 14, 2025;
January 2, 2015; January 22, 2026).
The formula $(3n)!/(3^n n!)$ in A052502 is credited to Ahmed Fares
(April 21, 2001). Accessed November 11, 2026.
Also A281183, \textsc{Formula}, Peter Bala, December 22, 2025,
\url{https://oeis.org/A281183}.
\bibitem{su}
Lien-Hung Su, \emph{Quadratic Derivations and Uniform Modular Periods
for Six Inverse-Function Sequences}, version 1, October 31, 2026.
Zenodo, \href{https://doi.org/10.5281/zenodo.23835577}{doi:10.5281/zenodo.23835577}.
Theorem 4.1 and Remark 4.2, pp.~7--8.
\bibitem{gessel}
Ira M. Gessel, \emph{Lagrange inversion}, Journal of Combinatorial Theory,
Series A \textbf{144} (2016), 212--249,
\href{https://doi.org/10.1016/j.jcta.2016.06.018}{doi:10.1016/j.jcta.2016.06.018}.
Author version: arXiv:1609.05988v1, Theorem 2.1.1,
\url{https://arxiv.org/abs/1609.05988v1}.
\end{thebibliography}
\bigskip\noindent\small
\copyright\ 2026 Haobo Ma and Wenlin Zhang.
Licensed under Creative Commons Attribution 4.0 International
(CC BY 4.0): \url{https://creativecommons.org/licenses/by/4.0/}.
\end{document}
